\section{Limitations}

This model, like any theoretical model, has omissions and limitations that must be acknowledged. This model's most important scope condition is that housing needs to be allocated via markets. The reviewed literature primarily focuses on North America and Western Europe, so the stylized facts explained by this model may be geographically limited. Even within these contexts, this model makes noteworthy assumptions. First, this model does not consider strategic action and exploitation \citep{desmondPoorPayMore2019}. It is noteworthy, however, that this model is nonetheless able to explain patterns of housing inequality and unaffordability. I am not suggesting that there is no exploitation in housing markets, but inequalities might persist without exploitative practices. Second, we know that land and housing are often used as assets and investment properties; therefore, housing supply depends on macroeconomic conditions \citep{bayerSpeculativeFeverInvestor2021}, which I have omitted from this model for simplicity. Macro-economic conditions could make it more or less profitable to invest in housing stock and, therefore, be a trigger for neighborhood change. Third, I also omitted geographical features and distance to amenities. I abstracted from exogenous amenities in my model to highlight the endogenous processes at work that sustain segregation and lead to changes in neighborhoods due to shocks. Although this provides many insights into these processes and can explain multiple stylized facts, it is fundamentally aspatial. It overlooks a substantial body of literature that documents the importance of amenities. These could further explain where groups segregate to \citep{gaigneWhoLivesWhere2022} and which neighborhoods change \citep{leeNaturalAmenitiesNeighbourhood2018}. Zoning has also been omitted, not just for simplicity but also due to endogeneity. Particularly, density zoning is a way to institutionalize areas where poor households cannot afford to move to; such laws likely stem from affluent households' preferences to keep the poor out of their areas \citep{rothwellDensityZoningClass2010}.

The process of moving is also highly stylized. In reality, searching for housing or new occupants takes time, is connected to transaction costs, depends on the availability of information, often involves intermediaries like real estate agents, and households might be discriminated against \citep{hanMicrostructureHousingMarkets2015}. Households typically move around life-course events, rather than marginally increasing their utility. Some sociologists have recently called for more realistic micro-foundations of theoretical models of residential stratification. Particularly, they review and present evidence that households first heuristically narrow down the neighborhood in which they search for a housing unit before assessing the qualities of available units more rationally \citep{krysanCycleSegregationSocial2017, bruchChoiceSetFormation2019}. However, the general assumption that households want to find a high-quality unit in a high-quality neighborhood but are constrained by their budget remains plausible (although there is some heterogeneity in preferences and decision-making), and the less realistic depiction of the housing search process simplifies the model and keeps it in line with economic theory. There are two ways, however, in which the utility maximization assumption may contribute to substantial shortcomings of the model. First, when households consider neighborhood and housing quality at different stages of the search process, these goods are not substitutes. It is, however, hard to say if this leads to over- or underestimation of segregation. Second, utility maximization models cannot explain the established link between increasing income inequality and increasing income segregation as well as the high segregation of affluent households (see section on robustness checks). Additionally, that households move to marginally increase their residential utility results in an overestimation of voluntary residential mobility in the model.

I have also framed this model within the context of renting, even though a large share of the population owns their home. There are arguably many differences between housing and rental markets but the reason for this framing is to illustrate two different \textit{roles} in markets, supply and demand. In rental markets, these roles are held by different people but in markets for owner-occupied housing, these might be the same people at different stages of the life course or even simultaneously when selling an old home and buying a new one. While there has been some research that institutional landlords behave differently as they have a stronger drive for profit (e.g., \cite{gorbackImpactInstitutionalOwners2025}), the general problem of uncertainty and the goal to get a good price and avoid large losses should hold for non-institutional sellers of owner-occupied housing as well. Whether this approximation works sufficiently well or owner-occupiers behave more distinctly in a way that is relevant for segregation and neighborhood formation, is an empirical question about the scope of this theory. It might be the case that the intrinsic motivation to invest in ones own house for your own comfort and the longer periods of stay in owner-occupied housing make such neighborhoods more stable compared to neighborhoods with more rental housing.